\section{Follow-up: real decisions and the cost of switching models}\label{sec:followup}

A reviewer of this report raised two objections. First, the constructed cases may be too easy and too
synthetic: he would rather see fewer cases with higher fidelity, taken from real agent runs. Second, prompt
caching is the most likely pushback against routing: \emph{``for longer chats you have to be confident in
savings for the smaller model.''} We ran one follow-up study for each. This chapter reports both. Nothing in
\cref{sec:intro,sec:background,sec:method,sec:validate,sec:dev,sec:holdout,sec:failures,sec:latency} was
re-measured or changed; the follow-ups are separate runs with their own evidence and their own
preregistration.

\begin{keyidea}
Two new questions, two new kinds of evidence. \textbf{Real decisions:} do the judges still help when the
decisions come from real agent sessions instead of hand-written cases? \textbf{Switching cost:} once a
conversation is long, does moving it to a cheaper model still save money, given that each model keeps its
own prompt cache?
\end{keyidea}

\subsection{Why real traces, and how the cases were built}\label{sec:traces}

\paragraph{What a trace-derived case is.} During earlier benchmark sessions, the bundle's in-session judge
faced thousands of real \emph{read-shortcut} decisions: ``should the agent just read this file, or think?''
A \term{trace-derived case} is one such decision point, rebuilt exactly as the judge saw it. We kept only
steps that the in-session judge had \emph{sent to the host model} (it fell back), because only then do we
know what the host did next. Steps the judge automated were excluded (\TrFastRouted\ of them): their
``outcome'' is the judge's own pick, so labelling them would be circular.

\paragraph{Rebuilding and fingerprinting.} The request bytes on the wire were never logged, so each request
was rebuilt from the session transcript. A rebuilt case was eligible only if it reproduced five logged
state numbers plus the logged hashes of the candidate ids and their order; \TrRebuildFail\ decisions failed
and were dropped. This is a fingerprint match, not a capture. From \TrMined\ eligible, de-duplicated
decisions we drew \TrDrawn\ (seed \TrSeed), stratified by suite and ``hard first'' using outcome-side
features only (more candidates, line-window candidates, clipped observations), never whether a logged judge
had been wrong. Whole tasks were assigned to holdout or dev before anyone labelled anything.

\paragraph{Labels.} Two reviewers labelled every case blind: they saw the judge's view and the host's next
one or two actions, but not the outcome-derived label or any judge's answer. They agreed with each other
at Cohen's $\kappa = \TrKappaAB$, and with the outcome label at only $\kappa = \TrKappaOutcome$. The final
label is the reviewers' shared label: \TrAgreed\ cases matched the outcome label, \TrOverridden\ overrode
it, and the \TrDropped\ cases where the reviewers disagreed were dropped. The holdout has \TrN\ cases
(\TrReasonLabels\ labelled \code{reason}, \TrReadCases\ labelled with a read); the dev split has \TrDevN\
and is a screen.

\begin{keyidea}
\textbf{``What the host did next'' is not the same as ``what was enough.''} The host model might read a file
it did not need, or read a different, equally good file. A careful reviewer asks a stricter question: from
what the judge could see, was one prepared read clearly sufficient? The two disagree often (the low $\kappa$
above), so the reviewers' label is the primary one and the host's next read is kept as a secondary check.
\end{keyidea}

\paragraph{Replayed in each judge's own form.} Jev received the bundle's own System One request; Laya and
the Qwen3 models went through the bundle's own local backends. GPT-6 Luna and Sol have no native form and
got the benchmark's choice question, as did nimble 9B and the tev1 models, because Ollama's System One
endpoint rejects the bundle's object-valued criteria (\cref{sec:newfindings}). The ``form'' column of
\cref{tab:trace} records which.

\paragraph{Preregistration.} The case hash, endpoints and decision rules were committed before the holdout
ran, and the runner refuses a holdout run unless they match. Two primary endpoints: the wrong-automatic rate
on all \TrN\ cases, and the number of \emph{correct automatic reads} among the \TrReadCases\ cases where a
read was enough (the savings a judge actually delivers). A new rule asks whether the read shortcut is worth
running at all: a judge is ``useful here'' if it makes at least \TrUsefulMinReads\ correct automatic reads
with a wrong-automatic upper 95\,\% bound below \TrUsefulMaxUpperPct\,\%. One erratum: a
\TrSmokeCases-case smoke test of the adapters on dev cases ran about \TrSmokeMinutes\ minutes before the
preregistration commit, although the file says it was committed before any judge saw a trace case. No
holdout case was involved. API spend for the whole study was \$\TrSpendTotal.

\subsection{Results on real decisions \texorpdfstring{\prereg}{[preregistered]}}\label{sec:traceresults}

\begin{table}[htbp]
\centering\small
\caption{Real read-shortcut decisions, holdout \prereg: \TrN\ cases, \TrReps\ repetitions, the bundle's
read-shortcut gate. Per-case majorities over repetitions. ``Upper'' is the upper end of the 95\,\% Wilson
interval of the wrong-automatic rate. ``Correct reads'' counts correct automatic reads among the
\TrReadCases\ cases whose label is a read. Latency pools all valid requests. The last row is the reference
that never acts.}
\label{tab:trace}
\begin{tabular*}{\textwidth}{@{\extracolsep{\fill}}l l r r r r r r r r@{}}
\toprule
 & & \multicolumn{2}{c}{Accuracy} & \multicolumn{2}{c}{Wrong automatic} & & Correct & \multicolumn{2}{c}{Latency (ms)} \\
\cmidrule(lr){3-4}\cmidrule(lr){5-6}\cmidrule(lr){9-10}
Judge & form & correct & 95\,\% CI & count & upper & automatic & reads (of \TrReadCases) & p50 & p95 \\
\midrule
\textbf{Jev 1.13} & native & \TrJevAccK/\TrN & \TrJevAccCI & \TrJevWaK & \TrJevWaUpper\,\% & \TrJevCovK & \TrJevReads & \TrJevPfifty & \TrJevPninetyfive \\
nimble 9B & adapted & \TrNimbleAccK/\TrN & \TrNimbleAccCI & \TrNimbleWaK & \TrNimbleWaUpper\,\% & \TrNimbleCovK & \TrNimbleReads & \TrNimblePfifty & \TrNimblePninetyfive \\
GPT-6 Luna & adapted & \TrLunaAccK/\TrN & \TrLunaAccCI & \TrLunaWaK & \TrLunaWaUpper\,\% & \TrLunaCovK & \TrLunaReads & \TrLunaPfifty & \TrLunaPninetyfive \\
GPT-6.1 Sol & adapted & \TrSolAccK/\TrN & \TrSolAccCI & \TrSolWaK & \TrSolWaUpper\,\% & \TrSolCovK & \TrSolReads & \TrSolPfifty & \TrSolPninetyfive \\
Qwen3 8B & native & \TrQwenEightAccK/\TrN & \TrQwenEightAccCI & \TrQwenEightWaK & \TrQwenEightWaUpper\,\% & \TrQwenEightCovK & \TrQwenEightReads & \TrQwenEightPfifty & \TrQwenEightPninetyfive \\
Qwen3 0.6B & native & \TrQwenPtSixAccK/\TrN & \TrQwenPtSixAccCI & \TrQwenPtSixWaK & \TrQwenPtSixWaUpper\,\% & \TrQwenPtSixCovK & \TrQwenPtSixReads & \TrQwenPtSixPfifty & \TrQwenPtSixPninetyfive \\
tev1 4B & adapted & \TrTevFourAccK/\TrN & \TrTevFourAccCI & \TrTevFourWaK & \TrTevFourWaUpper\,\% & \TrTevFourCovK & \TrTevFourReads & \TrTevFourPfifty & \TrTevFourPninetyfive \\
Laya base & native & \TrLayaAccK/\TrN & \TrLayaAccCI & \TrLayaWaK & \TrLayaWaUpper\,\% & \TrLayaCovK & \TrLayaReads & \TrLayaPfifty & \TrLayaPninetyfive \\
Qwen3 4B & native & \TrQwenFourAccK/\TrN & \TrQwenFourAccCI & \TrQwenFourWaK & \TrQwenFourWaUpper\,\% & \TrQwenFourCovK & \TrQwenFourReads & \TrQwenFourPfifty & \TrQwenFourPninetyfive \\
tev1 0.8B & adapted & \TrTevPtEightAccK/\TrN & \TrTevPtEightAccCI & \TrTevPtEightWaK & \TrTevPtEightWaUpper\,\% & \TrTevPtEightCovK & \TrTevPtEightReads & \TrTevPtEightPfifty & \TrTevPtEightPninetyfive \\
\midrule
\emph{Always fall back} & --- & \TrAlwaysK/\TrN & \TrAlwaysCI & 0 & --- & 0 & 0 & --- & --- \\
\bottomrule
\end{tabular*}

\end{table}
\howtoreadtable{A useful judge sits high in the ``correct reads'' column and low in the ``wrong automatic''
column. Accuracy is shown, but it rewards doing nothing: \TrReasonLabels\ of the \TrN\ labels are
\code{reason}, so a judge that never acts scores \TrAlwaysK/\TrN.}

\paragraph{No judge met the usefulness bar.} \TrNUseful\ of the \NumBase\ judges qualified. \TrNGeMinReads\
judges took at least \TrUsefulMinReads\ correct reads, and every one of them also made at least one wrong
automatic read. With \TrN\ cases, an upper bound below \TrUsefulMaxUpperPct\,\% is only possible with zero
wrong reads, so the rule effectively asked for a judge that saves several steps and never misfires.

\paragraph{Jev captured the most savings, and also misfired.} Jev took \TrJevReads\ of the \TrReadCases\
sufficient reads automatically and made \TrJevWaK\ wrong automatic reads (\TrJevWaPct\,\% of cases, upper
bound \TrJevWaUpper\,\%). GPT-6 Luna took \TrLunaReads\ with \TrLunaWaK\ wrong; GPT-6.1 Sol \TrSolReads\
with \TrSolWaK\ wrong, but \TrSolTimeouts\ of its \TrSolAnswers\ answers missed the \TrTimeoutMs\,ms timeout
(its median per repetition was \TrSolPfiftyRangeS\,s) and it cost \TrSolJevCostX$\times$ as much as Jev.
Ignoring the timeout, the same Sol answers would have taken \TrSolNoTimeoutReads\ reads with
\TrSolNoTimeoutWa\ wrong; that is hypothetical, because a slower judge also makes the host wait.

\paragraph{Jev, Luna and Sol are within noise.} Jev versus Luna: Holm-adjusted $p = \TrJevLunaAccPHolm$ on
accuracy and \TrJevLunaWaPHolm\ on wrong automatic reads; Luna versus Sol: \TrLunaSolAccPHolm\ and
\TrLunaSolWaPHolm. Jev versus Sol was missing from the run's contrast list; a post-hoc exact McNemar test
gives $p = \TrJevSolAccP$ (accuracy) and \TrJevSolWaP\ (wrong automatic). On the \TrReadCases\ read cases
alone, Jev's extra reads over Luna and Sol give $p = \TrJevLunaReadsP$ and \TrJevSolReadsP. No candidate
replaced Jev under the default-judge rule (\TrNReplaces\ did).

\paragraph{The local judges fail in two ways.} Laya and tev1 0.8B never cleared the \MinP\ gate, so they
never acted. Qwen3 4B never gave an answer at all (\cref{sec:newfindings}): all \TrQwenFourAbstained\ of
its answers became \code{reason} at probability 1. Nimble 9B, Qwen3 0.6B and Qwen3 8B acted freely and were
often wrong (\TrNimbleWaK, \TrQwenPtSixWaK\ and \TrQwenEightWaK\ wrong automatic reads).

\begin{figure}[htbp]
\centering
% Real decisions (trace holdout): correct automatic reads (of the read-labelled cases) against wrong
% automatic reads (of all cases), 95% Wilson whiskers on both axes. Labels: generated/labels/reads-holdout.tex.
\begin{tikzpicture}
\begin{axis}[benchmark,
  scale only axis, width=\ScatterAxisW, height=\ScatterAxisH,
  xmin=-1, xmax=14, ymin=-1.5, ymax=22,
  xtick={0,2,4,6,8,10,12}, ytick={0,5,10,15,20},
  xlabel={Correct automatic reads (of \TrReadCases\ cases where a read was enough)},
  ylabel={Wrong automatic reads (of \TrN)},
  xmajorgrids, ymajorgrids,
  legend style={at={(0.5,1)}, anchor=south, yshift=5pt, legend columns=4,
    /tikz/every even column/.append style={column sep=8pt}},
  legend cell align=left,
  clip=false,
]
\fill[okgreen!14] (axis cs:\ReadsZoneXa,\ReadsZoneYa) rectangle (axis cs:\ReadsZoneXb,\ReadsZoneYb);
\addlegendimage{area legend, fill=okgreen!14, draw=none}
\addlegendentry{rule-2 ``useful'' zone}
\addplot[only marks, mark=*, mark size=2.6pt, color=okblue,
  error bars/.cd, x dir=both, x explicit, y dir=both, y explicit, error bar style={okblue!55, line width=0.6pt}]
  table[x=reads, y=wa, x error minus=readsminus, x error plus=readsplus, y error minus=waminus, y error plus=waplus]
  {generated/data/reads-cloud.dat};
\addlegendentry{Cloud LLM}
\addplot[only marks, mark=square*, mark size=2.6pt, color=okorange,
  error bars/.cd, x dir=both, x explicit, y dir=both, y explicit, error bar style={okorange!55, line width=0.6pt}]
  table[x=reads, y=wa, x error minus=readsminus, x error plus=readsplus, y error minus=waminus, y error plus=waplus]
  {generated/data/reads-jev.dat};
\addlegendentry{Jev 1.13}
\addplot[only marks, mark=triangle*, mark size=3.0pt, color=okgreen,
  error bars/.cd, x dir=both, x explicit, y dir=both, y explicit, error bar style={okgreen!55, line width=0.6pt}]
  table[x=reads, y=wa, x error minus=readsminus, x error plus=readsplus, y error minus=waminus, y error plus=waplus]
  {generated/data/reads-local.dat};
\addlegendentry{Local model}
% Generated by build_assets.py: labels for panel reads-holdout. Do not edit.
\node[scatterlabel, anchor=south, scatterlabelmulti] at ($(axis cs:0.0000,0.0000)+(0.000pt,4.700pt)$) {tev1 0.8B,\\Qwen3 4B,\\Laya base};
\node[scatterlabel, anchor=west] at ($(axis cs:4.0000,3.0000)+(4.700pt,0.000pt)$) {tev1 4B};
\node[scatterlabel, anchor=west] at ($(axis cs:6.0000,9.0000)+(4.700pt,0.000pt)$) {Qwen3 0.6B};
\node[scatterlabel, anchor=west] at ($(axis cs:7.0000,2.0000)+(4.700pt,0.000pt)$) {GPT-6.1 Sol};
\node[scatterlabel, anchor=west] at ($(axis cs:7.0000,13.0000)+(4.700pt,0.000pt)$) {Qwen3 8B};
\node[scatterlabel, anchor=east] at ($(axis cs:8.0000,5.0000)+(-4.700pt,0.000pt)$) {GPT-6 Luna};
\node[scatterlabel, anchor=west] at ($(axis cs:9.0000,9.0000)+(4.700pt,0.000pt)$) {nimble 9B};
\node[scatterlabel, anchor=west] at ($(axis cs:10.0000,5.0000)+(4.700pt,0.000pt)$) {Jev 1.13};

\end{axis}
\end{tikzpicture}

\caption{Real decisions: correct automatic reads against wrong automatic reads, holdout \prereg, with
95\,\% Wilson intervals on both axes. Judges with identical results share one marker and one label. The
shaded strip is where rule 2 would call a judge useful (at least \TrUsefulMinReads\ correct reads and no
wrong ones).}
\label{fig:reads}
\howtoread{Right is more savings; up is more misfires. The ideal judge sits far right on the floor, inside
the shaded strip. No judge does: everyone that reaches the strip's left edge has also risen off the floor.
Long whiskers are honest: the horizontal axis has only \TrReadCases\ chances to succeed.}
\end{figure}

\paragraph{Why accuracy misleads here.} Always falling back scores \TrAlwaysK/\TrN, more than any judge
that acts. But \TrAlwaysSThreeK\ of those \TrAlwaysK\ come from the \TrSThree\ SWE-bench cases, where the
judge's view had been clipped before the issue text (\cref{sec:caveats}); on the other \TrSOneTwo\ cases
it scores \TrAlwaysSOneTwoK. The useful question is not ``how often is the judge right'' but ``how many
sufficient reads does it capture for each unnecessary one.''

\paragraph{What a wrong automatic read costs.} Every candidate here is a read-only file read. A wrong
automatic read replaces the host's next step with a read it did not need: one lost step and some extra
context, and no irreversible effect. That is a very different failure from clicking \emph{Buy now}
(\cref{sec:worked}). Its cost was assumed, not measured. Real states are also long, about
\TrJevInputTokens\ input tokens, so Jev cost \$\TrJevCost\ per million decisions here, \TrJevCostXFirst$\times$
its cost on the constructed holdout.

\paragraph{Against the host's own next read.} Scored against what the host actually read next (\TrHostReadCases\
cases where it read something), Jev has \TrJevHostReads\ correct automatic reads and \TrJevHostWa\ wrong;
Luna \TrLunaHostReads\ and \TrLunaHostWa; Sol \TrSolHostReads\ and \TrSolHostWa. None of Jev's \TrJevWaK\
wrong automatic reads matched the host's next read (\TrJevWaMatchHost). This analysis was chosen after the
results were seen and is not preregistered.

\subsection{Read this before generalizing}\label{sec:caveats}

\begin{caution}[title={Read this before generalizing}]
\begin{itemize}[leftmargin=1.1em]
  \item \textbf{Scope and selection.} These are \TrN\ decisions that some in-session judge had already
    \emph{declined} and sent to the host. They are not a sample of all decisions: about half are labelled
    ``a read was enough'' by design, against about \TrPoolReadPct\,\% of the full pool. Jev made
    \TrJevOwnWa\ wrong automatic reads on the \TrJevOwnN\ cases from its own sessions (cases its
    in-session self had declined) and \TrJevOtherWa\ on the \TrJevOtherN\ cases from other judges'
    sessions.
  \item \textbf{Label overrides.} Reviewers overrode the host-derived label in \TrOverridden\ of
    \TrKept\ kept cases, leaning strict. Of the \TrHoldOverrides\ holdout overrides, \TrOverReadToReason\
    turned a read into \code{reason}, \TrOverReasonToRead\ turned \code{reason} into a read and
    \TrOverReadToRead\ swapped one read for another. Reviewers saw the host's next actions; they were blind
    only to the proposed label and to judge answers.
  \item \textbf{Task clipping.} Every case ran at a \TrStateBudget-character state budget (the bundle
    default is \TrStateDefault). In the \TrSThree\ SWE-bench holdout cases the task text the judge saw ended
    after about \TrTaskSeen\ characters, inside the operating rules: the judge never saw the issue.
  \item \textbf{Qwen3 4B is an artifact, not caution.} Its zero wrong automatic reads come from a defect
    in the candidate path, not from the model's judgement.
  \item \textbf{Clustering and source.} The \TrN\ cases come from \TrGroups\ task groups; Jev's
    \TrJevWaK\ wrong reads fall in \TrJevWrongGroups. All cases come from this bundle's own benchmark
    sessions, not human-driven traffic, and the reviewers are Anthropic models.
  \item \textbf{Erratum.} A \TrSmokeCases-case dev smoke test ran before the preregistration commit
    (\cref{sec:traces}); no holdout case was involved.
\end{itemize}
\end{caution}

\subsection{Caching: does switching models erase the saving?}\label{sec:caching}

\paragraph{The question.} Routing sends some turns of a conversation to a cheaper model. Providers keep a
prompt cache \emph{per model}: re-sending a long conversation to the same model is cheap, because most of it
is read from cache. After a switch, the new model must first \emph{write} the whole conversation into its
own cache, and cache writes are the most expensive input tokens.

\begin{keyidea}
Three prices matter. On the Opus host, a cached read costs \$\CaOpusRead\ per million tokens, a cache write
\$\CaOpusWrite\ and an uncached input token \$\CaOpusIn. A model switch turns cheap reads into expensive
writes once (a \term{cache rebuild}). A cheaper model is not automatically cheaper in a long chat: Sonnet
reads its cache at \CaReadPriceX$\times$ Opus's price, while every other token class costs \CaOtherPriceX$\times$.
\end{keyidea}

\paragraph{The data.} This study made no model calls: it re-read \CaRuns\ earlier benchmark runs and their
session logs and formed \CaPairs\ matched pairs (\CaRoutedPairs\ of them routed): a routed run and the plain
run on the same host, task, campaign and repetition. Ratios are geometric means of routed cost over plain
cost; intervals bootstrap over tasks (\CaBoot\ resamples). The cost field is a price table, not billing:
it reproduces list prices exactly.

\begin{table}[htbp]
\centering\small
\caption{Price table behind every cost in this section (US dollars per million tokens).}
\label{tab:prices}
\begin{tabular*}{0.75\textwidth}{@{\extracolsep{\fill}}l r r r r@{}}
\toprule
Model & uncached input & cache read & cache write & output \\
\midrule
Fable 5.1 & \CaFableIn & \CaFableRead & \CaFableWrite & \CaFableOut \\
Opus 5.5 & \CaOpusIn & \CaOpusRead & \CaOpusWrite & \CaOpusOut \\
Sonnet 5 & \CaSonnetIn & \CaSonnetRead & \CaSonnetWrite & \CaSonnetOut \\
\bottomrule
\end{tabular*}
\end{table}

\paragraph{The same-cell contrast.} The fair comparison holds the routing configuration fixed and changes
only the session length. \Cref{tab:samecell} does that for single-turn and \CaTurns-turn sessions.

\begin{table}[htbp]
\centering\small
\caption{Cost ratio (routed or swapped run over plain run on the same host; below 1 saves money) for the
same cell in single-turn and \CaTurns-turn sessions, with 95\,\% bootstrap intervals. The last row never
switches models within a session: it is the control.}
\label{tab:samecell}
\begin{tabular*}{\textwidth}{@{\extracolsep{\fill}}l r r r r@{}}
\toprule
 & \multicolumn{2}{c}{Single-turn} & \multicolumn{2}{c}{\CaTurns-turn} \\
\cmidrule(lr){2-3}\cmidrule(l){4-5}
Cell (host) & cost ratio & pairs & cost ratio & pairs \\
\midrule
orch-default-opus vs plain-opus (Opus) & \CaOpusRoutedOne\ [\CaOpusRoutedOneCI] & \CaOpusRoutedOnePairs & \CaOpusRoutedFour\ [\CaOpusRoutedFourCI] & \CaOpusRoutedFourPairs \\
orch-default vs plain (Fable) & \CaFableRoutedOne\ [\CaFableRoutedOneCI] & \CaFableRoutedOnePairs & \CaFableRoutedFour\ [\CaFableRoutedFourCI] & \CaFableRoutedFourPairs \\
plain-sonnet vs plain-opus (no switch) & \CaControlOne\ [\CaControlOneCI] & \CaControlOnePairs & \CaControlFour\ [\CaControlFourCI] & \CaControlFourPairs \\
\bottomrule
\end{tabular*}
\end{table}
\howtoreadtable{Read each row left to right. On the Opus host, routing went from a small saving
(\CaOpusRoutedOne$\times$) to a loss (\CaOpusRoutedFour$\times$). The control, which simply runs Sonnet
instead of Opus and never switches, moved almost the same way (\CaControlOne$\times$ to
\CaControlFour$\times$). So most of the Opus-host loss is not the switch itself. On the Fable host routing
still saved money in longer sessions, though less.}

\paragraph{Price and volume, about half each.} Re-pricing the control's \emph{identical} tokens at Sonnet's
prices gives \CaPriceOne$\times$ Opus in single-turn sessions, where cache writes dominate, but
\CaPriceFour$\times$ in \CaTurns-turn sessions, where cache reads are \CaReadShareFour\,\% of Opus's cost
(against \CaReadShareOne\,\% single-turn). Sonnet also used more: \CaRequestsFour$\times$ the requests and
\CaCacheReadFour$\times$ the cache-read tokens in \CaTurns-turn sessions. Of the rise from \CaControlOne$\times$
to \CaControlFour$\times$, about \CaPriceShare\,\% is the price edge eroding and \CaVolumeShare\,\% is extra
volume (volume factor \CaVolumeOne$\times$ to \CaVolumeFour$\times$).

\paragraph{How much the rebuilds cost.} There were \CaSwitches\ model switches: \CaSwitchBoundary\ at a turn
boundary and \CaSwitchMid\ in the middle of a turn. Turn 2 is where most routes move to the host model
(\CaTurnTwoSwitches\ switches); on the Opus host turn 2 cost \CaOpusTurnTwo$\times$ the plain run's turn 2
(\CaOpusTurnTwoEx$\times$ without the rebuild), and on Fable \CaFableTurnTwo$\times$ (\CaFableTurnTwoEx$\times$).
The share of routed cost spent on rebuilds depends strongly on the kind of session:

\begin{table}[htbp]
\centering\small
\caption{Rebuild cost after model switches, by stratum (share of the routed runs' cost).}
\label{tab:rebuild}
\begin{tabular*}{\textwidth}{@{\extracolsep{\fill}}l r r r@{}}
\toprule
Stratum & pairs & sessions with a switch & rebuild share \\
\midrule
Single-turn routing & \CaRbOnePairs & \CaRbOneSwitch & \CaRbOneShare\,\% \\
\CaTurns-turn, Opus host & \CaRbFourOpusPairs & \CaRbFourOpusSwitch & \CaRbFourOpusShare\,\% \\
\CaTurns-turn, Fable host & \CaRbFourFablePairs & \CaRbFourFableSwitch & \CaRbFourFableShare\,\% \\
SWE-bench, mid-turn escalation & \CaRbPrimaryPairs & \CaRbPrimarySwitch & \CaRbPrimaryShare\,\% (\CaRbPrimaryEffortShare\,\%$^{a}$) \\
SWE-bench, turn-start routers & \CaRbRoutersPairs & \CaRbRoutersSwitch & \CaRbRoutersShare\,\% \\
\bottomrule
\end{tabular*}
\par\smallskip\raggedright\footnotesize $^{a}$ Including cache rewrites caused by effort-setting changes on
the same model.
\end{table}

The shipped default router made no mid-turn switches on SWE-bench (\CaRulesSThreeSwitches\ switches) and
came in at \CaRulesSThree$\times$. Pooled ``rebuild share'' totals across strata are misleading: several
routed cells share one plain anchor run (\CaUniqueAnchors\ anchors behind \CaRoutedPairs\ pairs).

\paragraph{Warm first requests.} Runs ran concurrently, so some sessions found their system prompt
already cached by another run, more often for routed runs (\CaWarmOpusTreat\,\% against
\CaWarmOpusAnchor\,\% of anchors on the Opus host, \CaTurns-turn). Re-normalizing both arms to all-cold or
all-warm moves the Opus routing ratio within \CaWarmOpusRange\ (raw \CaWarmOpusRaw), Fable within
\CaWarmFableRange\ (raw \CaWarmFableRaw) and the control within \CaWarmControlRange\ (raw
\CaWarmControlRaw). The conclusions hold across these ranges.

\begin{caution}[title={What the caching study does not measure}]
\begin{itemize}[leftmargin=1.1em]
  \item \textbf{Longer and real chats.} Only \CaScenarios\ synthetic scenarios, at most \CaTurns\ turns,
    about \CaTurnGapS\,s apart, and no human chats. \CaOpusFourDevPairs\ of the \CaOpusFourPairs\ Opus
    \CaTurns-turn pairs (\CaFableFourDevPairs\ of \CaFableFourPairs\ on Fable) come from the
    \CaDevScenarios\ development scenarios, which were used for tuning.
  \item \textbf{Cache expiry.} The cache expires after about \CaExpiryMin\ minutes, and turns were seconds
    apart, so the plain run never paid for an expired cache. With real think time both sides would pay
    rewrites.
  \item \textbf{The switch in isolation.} Routed cells also change effort settings and add the
    orchestrator, and no arm ``decides once and never switches.''
  \item \textbf{Receipts and billing.} \CaReceipts\ of the \CaRoutedPairs\ routed pairs carry efficiency
    receipts, and the costs are list-price estimates, not invoices.
\end{itemize}
\end{caution}

\paragraph{The proposed experiment (not run).} The caching study proposes a direct test: \CaPropScenarios\
new, preregistered, \CaPropTurns-turn scenarios with checks after each turn; Opus and Fable hosts; a new
``sticky'' arm that chooses a model once per session and never switches, to isolate the cost of switching;
a factor with \CaPropGapMin-minute gaps between turns so the cache expires; \CaPropReps\ repetitions; about
\$\CaPropUSD\ and \CaPropHours\ hours unattended. \textbf{It needs approval and has not been run}; nothing in
this report depends on it.

\subsection{New findings for the bundle}\label{sec:newfindings}

Building the trace cases surfaced four defects in the bundle. Each is documented with evidence in
\path{evals/judge_bench/traces/}. \textbf{None was fixed in this work}, and none changes a number above
except where stated.

\begin{enumerate}
  \item \textbf{The state budget clips the task.} The read-shortcut state keeps the \emph{beginning} of the
    task message. At the \TrStateBudget-character budget used in these sessions, the SWE-bench task text
    ended after about \TrTaskSeen\ characters, inside the operating rules, so the judge never saw the issue.
    The bundle's default budget is \TrStateDefault; the shorter budget came from the benchmark profiles.
  \item \textbf{The escalation and phase judges see the wrong task text.} Their state takes the first user
    message as the task, and in Amplifier sessions that message is the injected system-reminder block.
    Neither judge ran in any logged session, so no measured result is affected.
  \item \textbf{Ollama's System One endpoint rejects object-valued criteria.} A Jev-style backend pointed at
    Ollama's endpoint cannot answer read-shortcut questions (HTTP 400). This is why nimble 9B and the tev1
    models ran in the adapted form here.
  \item \textbf{Qwen3 4B never answers on the candidate path.} Its first generated token is prose
    (``We\dots''), so no option letter receives probability and everything goes to \code{reason}. The
    bundle's question path has a guard and a retry for this; the candidate path has neither. Qwen3 4B's
    zero wrong automatic reads in \cref{tab:trace} are this artifact.
\end{enumerate}
Escalation, phase and tool-risk decisions could not be turned into trace cases at all: their payloads were
never logged live, or carry too little to support an honest label.
